\subsection{Comparison with the base paper}
\label{sec:audit}

This subsection gathers, in one place, how our results line up with the
base paper \cite{Koenig_2024}. We first compare our recipe with the
paper's stated methodology field by field
(Table~\ref{tab:audit}), so that no difference is silently absorbed into
``the paper trained longer''. The table uses the settings recorded in our
own planning notes and code, and marks as open every field those do not
determine. One such
field, the MLP baseline's optimizer settings, we later resolved by reading
the paper's own published code, which is how the follow-up runs of
Section~\ref{sec:phase4} came about.

\begin{table*}[t]
\centering
\footnotesize
\renewcommand{\arraystretch}{1.15}
\caption{Field-by-field comparison with the base paper \cite{Koenig_2024}.}
\label{tab:audit}
\begin{tabularx}{\textwidth}{@{}L{2.7cm} L{4.2cm} L{4.2cm} Y@{}}
\toprule
\hdrrow \thh{Field} & \thh{Base paper} & \thh{Ours} & \thh{Verdict}\\
\midrule
KAN architecture & $[2,10,2]$, $G=5$, 240 parameters & identical & Matches.\\
Edge basis & Gaussian RBF, $\exp(-r^2/2h^2)$ & identical, after fixing a missing factor $\tfrac12$ & Matches. The fix predates every reported run.\\
MLP baseline & $[2,50,2]$, tanh, 252 parameters & identical, plus a working $[2,14,8,8,2]$ SiLU variant & Matches literally, and trains under the paper's own optimizer settings (below).\\
Integrator & adaptive Tsit5 (Julia) & fixed-step Tsit5, $h=0.05$ & Explained: the fixed step's own error is $4.7\times10^{-14}$ in MSE (Table~\ref{tab:solver-ablation}).\\
Gradients & continuous adjoint & direct autograd & Explained: they differ in memory and time, not in the loss minimised (Section~\ref{sec:track-a}).\\
Framework & compiled Julia / SciML & interpreted PyTorch, CPU & Explained: wall-clock is not comparable across the two (Section~\ref{sec:budget}).\\
Epoch budget & $10^4$ (KAN), $10^5$ (MLP) & $10^4$, and $5\times10^4$ in Section~\ref{sec:phase4} & Explained: Section~\ref{sec:phase4}.\\
Optimizer, learning rate (MLP baseline) & Adam, $10^{-2}$, Glorot init $\div\,10^5$, found in the paper's own published code & Adam, $2\times10^{-3}$ by default, Glorot gain $1.0$, clip $1.0$ & \textbf{Resolved}: this mismatch, not the architecture, caused the earlier training failure. The paper's own KAN learning rate remains unconfirmed.\\
Loss normalisation & not recorded & mean over points and components & \textbf{Open}; affects absolute loss values only.\\
\bottomrule
\end{tabularx}
\end{table*}

\paragraph{The efficiency claim.} The paper reports its KAN-ODE reaching a
training loss of $2.6\times10^{-5}$ in $10^4$ epochs and its MLP needing
$10^5$ epochs for $3.0\times10^{-5}$. Our extended runs locate both targets
directly. Our KAN-ODE reaches $2.6\times10^{-5}$ at epoch $13{,}632$,
within $1.4\times$ of the paper's KAN. Our SiLU MLP reaches
$3.0\times10^{-5}$ at epoch $8{,}585$, nearly $12\times$ sooner than the
paper's MLP. The literal $[2,50,2]$ tanh MLP never gets below $1.0$ under
our default optimizer settings, but under the paper's own learning rate and
initialization it crosses $10^{-4}$ at epoch $39{,}695$ and had not reached
$3.0\times10^{-5}$ by epoch $50{,}000$: a trend consistent with, but short
of confirming, the paper's own $10^5$-epoch target. Allowing for the open
loss-normalisation row, the KAN side of the efficiency claim is reproduced.
The tenfold advantage over the literal MLP baseline, however, depends on
that baseline's optimizer settings as much as on its architecture, since the
identical network trains once those settings are corrected.

\paragraph{The accuracy claim.} At the paper's own $10^4$-epoch budget we
reproduce the extrapolation advantage: KAN-ODE is $5.4\times$ better than a
parameter-matched MLP (Section~\ref{sec:kan-vs-mlp}), and the gap widens to
$40.6\times$ at twice the horizon (Section~\ref{sec:crossdomain}). At five
times that budget the ranking reverses and the MLP is $6.8\times$ better
(Section~\ref{sec:phase4}), so we read the paper's accuracy claim as a
statement about moderate training budgets rather than about the two
function classes at convergence.

\paragraph{The three fixed design choices.} None of the paper's three fixed
choices turns out to be load-bearing on this benchmark. Tsit5 is not needed
for accuracy: above order two, every solver lands within $3.8\%$ of the
others, and Midpoint matches Tsit5 at $29\%$ of the cost
(Section~\ref{sec:solver-results}). The switch from B-spline to RBF is
justified on cost rather than accuracy: the two are statistically tied at
both budgets, with RBF $3.9\times$ cheaper per epoch
(Section~\ref{sec:basis-results}). And the adjoint method's memory saving,
while real at every trajectory length we tested, never matters at our
scale, while it costs $1.9$--$2.1\times$ in wall-clock
(Section~\ref{sec:track-a}).

\paragraph{Wall-clock.} Our absolute training times are not comparable with
the paper's. Our integration loop runs uncompiled in PyTorch on CPU,
whereas the paper compiles its loop in Julia's SciML ecosystem
(Table~\ref{tab:cost}). All our timing comparisons are therefore made only
between our own runs, on the same hardware.

\begin{keypoint}[Outcome of the comparison]
The $10^4$-epoch budget gives trustworthy \emph{rankings} for the solver and
basis ablations: the RBF/B-spline tie and Euler's last place both survive
five times the budget. It does not for the KAN-versus-MLP comparison, which
reverses: the paper's extrapolation advantage is an early-training effect
under our protocol, and its speed advantage over the literal MLP baseline
depended on that baseline's default optimizer settings, not on its
architecture. Corrected, the identical network trains, though it still does
not close the gap with either KAN-ODE or the SiLU MLP within our budget.
Every result is single-seed.
\end{keypoint}
